\section{Sparse packet arrays}
\label{s:section}

The transverse saving gives a strict improvement when each individual
packet encounters little measure.  The broad--narrow organization follows
Bourgain--Guth \citep[Section~2]{BourgainGuth2011} and its fractal
Schr\"odinger implementation in Du--Zhang \citep[Section~3]{DuZhang2019}.
Here the iteration must also preserve the packet moment condition under
parabolic rescaling.  Broad leaves gain from
Proposition~\ref{t:transverse}; at the stopping scale the packet moment
condition improves the square-root-scale density in the refined fractal
estimate.

\begin{proposition}[Sparse estimate]\label{s:sparse}
There is $c_s>0$, depending only on the fixed dimension and packet
boxes, with the following property.  Fix $\tau>0$ and
$L^\tau\le D\le L$.  Let $U$ be a full-dimensional compact-profile
array of scale $L$, with coefficient energy
$\mathcal E=\sum_\alpha|c_\alpha|^2$.  Let $\mu$ be a nonnegative
measure on the fixed native time range, obeying
\[
 \mu(B(z,r))\le mr^n\quad(r\ge1),\qquad
 \int\left(1+\frac{|x-y_\alpha-2t\omega_\alpha|}{L}\right)^{-M}
          \dd\mu(x,t)\le\frac{mL^{n+1}}D
\]
for every nonzero packet label, as in Definition~\ref{d:measure}.
For sufficiently many bounded profile derivatives,
\begin{equation}\label{s:estimate}
 \sum_{q\text{ unit}}\mu(q)^\beta\sup_q|U|^2
       \le C_\tau m^\beta L^{2n/(n+1)}D^{-c_s}\mathcal E,
       \qquad \beta=\frac2{n+1}.
\end{equation}
The same estimate holds for any fixed finite set of field derivatives
if the original profiles have sufficiently many more derivatives.
The exponent $c_s$ is independent of $\tau$; the constant and the
required finite differentiability may depend on $\tau$.  Restrictions
to submeasures and fixed enlargements of the native time range are
allowed.
\end{proposition}

\subsection{The local functional and exact layers}

Choose a dyadic $K\asymp L^v$, with $v>0$ chosen at the end, and put
\[
 P=\frac2{1-\beta}=\frac{2(n+1)}{n-1}=q_{n-1},\qquad
 \mathcal T_\beta(V,\nu)
       =\sum_{Q\text{ a }K^2\text{-cube}}
                      \nu(Q)^\beta\norm{V}_{L^P(Q)}^2.
\]
This is the functional from Section~\ref{d:section}; its square root
is a seminorm.  For a fixed finite set of multiindices $a$, unit-cube
Sobolev followed by H\"older gives
\begin{equation}\label{s:sobolev}
 \begin{split}
 \sum_{q\subset Q}\mu(q)^\beta\sup_q|U|^2
 &\lesssim\sum_a\sum_{q\subset Q}\mu(q)^\beta
                          \norm{\partial^aU}_{L^P(q)}^2\\
 &\le\sum_a\mu(Q)^\beta\norm{\partial^aU}_{L^P(Q)}^2.
 \end{split}
\end{equation}
Here $2/P=1-\beta$, and the grids may be chosen nested.  The cube
Sobolev inequality controls closures, so boundary suprema cause no
problem.  Differentiating a compact atom produces finitely many atoms
with the same labels and comparable energy: carrier frequencies are
bounded, and native derivatives carry only bounded powers of $L^{-1}$.
It suffices to estimate $\mathcal T_\beta$ for these arrays.

First restrict space to boxes of side $L^2$.  On the stipulated time
range a compact packet meets only boundedly many fixed enlargements of
these boxes, because its velocity is bounded.  The retained coefficient
energies therefore sum with bounded overlap.

Apply Lemma~\ref{d:layers} to each localized compact array.  Its exact
Fourier layers are indexed by $(k,u)\in\Z^n\times\R$, with an integrable
common weight $w(k,u)$ and exterior phase $e^{iut/L^2}$.  After division
by that weight, every subarray $A$ has exact input squared norm at most
$C\sum_{\alpha\in A}|c_\alpha|^2$.  Its frequency support is within
$O(L^{-1})$ of $\omega_\alpha+k/L$, and on a fixed slightly enlarged
native time range its packets have any prescribed finite spatial decay
order about the original trajectories and bounded native frequency
seminorms.  These properties follow from the native Fourier identity
$u=\zeta+|\eta|^2$ and the summable phase-space Gram rows; they apply
to every subarray with the same constant.

Minkowski for $\mathcal T_\beta^{1/2}$ returns a bound for the original
array after integrating the layer weights.  The exterior phase has
modulus one and is removed before the iteration.  For the derivative
version, differentiate compact atoms before taking layers; no derivative
of this exterior phase enters the iteration.
Choose a small fixed $\gamma>0$ later.  Layers with $|k|>L^\gamma$
have arbitrarily small negative-power contribution by choosing the
weight decay last.  Indeed, after normalizing $m=\mathcal E=1$ on one
original spatial box, the relevant masses, cell counts, and field
suprema are polynomial in $L$.  The layer input bound and bounded
frequency volume give the uniform supremum bound needed for this tail
estimate.  Henceforth one retained normalized layer is fixed.

\subsection{A fixed gain on broad cubes}

At packet scale $l$, group labels into $1/K$-frequency bins and write
$V=\sum_\theta V_\theta$.  Every subfield is exact and has support
within $O(l^{-1})$ of its label bin plus the common shift $k/l$.
Our choices will guarantee $l\gg K^2$ and $|k|/l=o(1)$.
Call a bin active on $Q$ when
\begin{equation}\label{s:activity}
 \norm{V_\theta}_{L^P(Q)}\ge K^{-A}\norm V_{L^P(Q)},
\end{equation}
where $A>n+1$ is fixed sufficiently large.  The $O(K^n)$ inactive
bins have total norm at most $CK^{n-A}\norm V_{L^P(Q)}$, which can
be absorbed.

\begin{lemma}[The geometric alternative]\label{s:simplex}
For a fixed sufficiently large $C_0$, either all active centers lie
within $C_0/K$ of an affine hyperplane, or $n+1$ active bins have
direction determinant at least $cK^{-n}$ throughout their exact
supports.  A common frequency shift does not affect this alternative.
\end{lemma}

\begin{proof}
Choose one center $\xi_0$, and then greedily choose a center of largest
distance from the affine span of the previous choices.  If a maximum
distance is at most $C_0/K$, all centers lie in that neighborhood of
the current affine span, hence of a hyperplane.  Otherwise the $n$
Gram--Schmidt heights $d_1\ge\cdots\ge d_n$ all exceed $C_0/K$.
In the resulting orthonormal basis, the column matrix of
$\xi_j-\xi_0$ is upper triangular with diagonal $d_j$ and every
entry in row $j$ bounded by $d_j$, by the greedy choice.  Factoring
out these row heights leaves a unit upper-triangular matrix with
bounded entries and bounded inverse, with constants depending only on
$n$.  Thus the inverse original matrix has norm $O_n(d_n^{-1})$.
Perturbing the vertices by $O(K^{-1})$ changes its determinant by a
relative $O_n((Kd_n)^{-1})$, which is small for sufficiently large
$C_0$.  Its determinant is at least $\prod_jd_j\gtrsim K^{-n}$.
The exact support enlargements are smaller than these bin perturbations.
The determinant of the spacetime directions $(2\xi_j,1)$ is a fixed
multiple of this simplex determinant.  Common translation preserves it;
unit normalization changes it only by bounded factors because the
frequencies remain bounded.
\end{proof}

Call the first case narrow and the second broad.

\begin{lemma}[Broad gain]\label{s:broad}
There exists $c_b>0$, depending only on the fixed frequency box and
$n$, such that if $K\le l^{c_b}$, then in every box of side $O(l^2)$,
\begin{equation}\label{s:broad-bound}
 \sum_{Q\text{ broad}}\nu(Q)^\beta\norm V_{L^P(Q)}^2
       \lesssim m_0^\beta l^{2n/(n+1)-c_b}\mathcal E_0.
\end{equation}
Here $\nu(B(z,r))\le m_0r^n$ for $r\ge1$, and the input squared
norms of the total field and every subfield are at most $C\mathcal E_0$.
The fixed $C$ may affect the implicit constant, but the gain is
independent of packet differentiability and decay orders.
\end{lemma}

\begin{proof}
If no positive $c_b$ works, normalize $m_0=\mathcal E_0=1$ and the
fixed input norm bound to obtain a sequence with $l\to\infty$,
$K=l^{o(1)}$, and broad contribution at least
$l^{2n/(n+1)-o(1)}$.  Polynomially small mass and norm levels are
negligible by the polynomial number of cubes and uniform exact-field
supremum bound.  Pigeonhole a common active transverse tuple and dyadic
levels $\nu(Q)\asymp\chi$, $\norm V_{L^P(Q)}\asymp A_0$ on $N$
cubes.  Tuple choices cost $K^{O_n(1)}$ and level choices cost logarithms,
all subpower here.  We have
\begin{equation}\label{s:broad-saturation}
 \chi^\beta A_0^2N\ge l^{2n/(n+1)-o(1)},\qquad
 \chi\le l^{o(1)},\qquad \chi N\lesssim l^{2n}.
\end{equation}

Each chosen active field has supremum at least
$K^{-A-(n-1)}A_0$ on some unit cube in $Q$, since
$|Q|^{1/P}=K^{n-1}$.  Pigeonholing the $n+1$ relative integer offsets
of these witnesses costs $K^{O_n(1)}$.  Translate the fields separately
by these fixed offsets so that their witnesses lie on a common unit
cube associated to each retained $Q$.  These cubes are distinct;
translations modulate exact inputs and preserve supports and norms.
Lemma~\ref{t:multilinear} applies with determinant
$K^{-n}=l^{-o(1)}$, and gives
\begin{equation}\label{s:product-count}
 A_0^{2(n+1)/n}N\le l^{o(1)}.
\end{equation}
In particular the inverse-transversality constant is controlled; no
unspecified dependence is used.

For completeness, take limiting base-$l$ logarithms $(x,y,z)$ of
$(\chi,A_0,N)$ along a subsequence.  The retained levels give polynomial
bounds, so these subsequences exist.  Equations~\eqref{s:broad-saturation}
and~\eqref{s:product-count} imply
\[
 x\le0,\quad x+z\le2n,\quad \frac{2(n+1)}n y+z\le0,\quad
 \beta x+2y+z\ge\frac{2n}{n+1}.
\]
They force equality throughout
\[
 \frac{2n}{n+1}\le\beta x+2y+z
 \le\frac{2x+z}{n+1}
 \le\frac{2n+x}{n+1}\le\frac{2n}{n+1}.
\]
Thus
\begin{equation}\label{s:saturated-sizes}
 \chi=l^{o(1)},\qquad N=l^{2n+o(1)},\qquad
 A_0=l^{-e+o(1)}.
\end{equation}
The attached witness cubes have $l^{o(1)}r^n$ counts for $r\ge1$:
their full disjoint $K^2$-cubes, each of mass at least $l^{-o(1)}$,
lie in a ball of radius $r+CK^2$ whenever their witnesses lie in a
radius-$r$ ball.  This gives at most
$l^{o(1)}(r+CK^2)^n\le l^{o(1)}r^n$ witnesses.
The translated fields and \eqref{s:saturated-sizes} now contradict
Proposition~\ref{t:transverse}.  This proves a fixed positive gain;
decrease it if necessary so the same $c_b$ occurs in both the hypothesis
and conclusion of the lemma.
\end{proof}

\subsection{Narrow transport and the measure on each branch}

On a narrow cube apply the codimension-one decoupling inequality in
Lemma~\ref{d:strip-decoupling} to the full sum of active bins.  Absorb
the inactive sum using \eqref{s:activity}, and then extend the right
side to all bins, with their complete child fields.  With any prescribed
large $N_1$, the squared-norm bound is
\begin{equation}\label{s:narrow-decoupling}
 \norm V_{L^P(Q)}^2\lesssim_\eps K^\eps
       \sum_\theta\sum_{a\in\Z^{n+1}}(1+|a|)^{-N_1}
                   \norm{V_\theta}_{L^P(Q+K^2a)}^2.
\end{equation}
The decoupling hypotheses hold because the active labels are in an
$O(K^{-1})$ hyperplane strip, the exact support enlargement is
$O(l^{-1})\ll K^{-1}$, and after removing the normal linear phase the
normal quadratic contribution has vertical thickness $O(K^{-2})$.
The exponent is exactly $P=q_{n-1}$.  Smooth cutoff localization and
bounded overlap of projected tangential caps, as in
Lemma~\ref{d:strip-decoupling}, allow arbitrary prescribed active
subfields, including overlapping enlarged supports.

For a child bin centered at $v_0$, use
\begin{equation}\label{s:child-map}
 \Phi_{v_0}(x,t)=\left(\frac{x-2v_0t}{K},\frac t{K^2}\right),
 \qquad l'=l/K,\quad y_\alpha'=y_\alpha/K,\quad
 \omega_\alpha'=K(\omega_\alpha-v_0).
\end{equation}
After removing the carrier, $V_\theta=K^{-n/2}V_\theta'\circ\Phi_{v_0}$.
Native variables, coefficient energies, and phase-space counting remain
unchanged.  For the direct pushforward $\nu'=(\Phi_{v_0})_\#\nu$,
\begin{equation}\label{s:direct-measure}
 \nu'(B(z,r))\le C K^{n+1}m_0r^n,\qquad
 \int\left(1+\frac{|x'-y_\alpha'-2t'\omega_\alpha'|}{l'}\right)^{-M}
                  \dd\nu'\le\frac{m_0l^{n+1}}D.
\end{equation}
Indeed the inverse ball is covered by $O(K)$ balls of radius $O(Kr)$.
The normalized transverse distance in the moment integrand is unchanged
pointwise, and $K^{n+1}(l')^{n+1}=l^{n+1}$.

Lemma~\ref{d:transport} gives, for each fixed shift $a$, a
finite sum of child functionals against translates of this direct
pushforward, with total factor
\begin{equation}\label{s:transport-power}
 CK^{2-(n+2)\beta}=CK^{-\beta}.
\end{equation}
To see the exponent, the Jacobian is $K^{n+2}$, the normalized field
contributes $K^{-n}$ in squared norm, and $2/P=1-\beta$; the product
is $K^{(n+2)(1-\beta)-n}$.  Group disjoint shifted parent images by
the child $K^2$-cube containing the image of their centers.  H\"older
within the group bounds the sum by total pushed-forward mass to power
$\beta$ times the norm on a bounded enlargement.  Subdividing the
enlargement introduces only finitely many neighboring cube translates.
There is no factor for the number of parent cubes in a group.

We keep each shift and each finite neighbor choice as a separate
weighted branch of the iteration.  The shift coefficients in
\eqref{s:narrow-decoupling} are summable, uniformly in their truncation;
the finite neighbor choices cost a fixed constant.  Thus the total
weighted coefficient energy of the children is at most a fixed
constant times the parent energy.  The bins themselves partition the
packet labels.  This bookkeeping permits each branch to retain a
measure with parameter $C K^{n+1}m_0$, without aggregating all
translated measures into one.

Truncate $|a|\le L^\eta$ for a sufficiently small fixed $\eta>0$.
The omitted tail is negligible by choosing $N_1$ last: exact fields
have polynomial global supremum bounds, and normalized relevant
masses, cells, and depths are polynomially bounded.  Each remaining
child translate has displacement $d$ with $|d|\lesssim K^2L^\eta$.
Its effect on the child's transverse displacement is
$d_x-2\omega_\alpha'd_t$.  The child label velocities are bounded, and
we impose
\begin{equation}\label{s:shift-smallness}
 K^2L^\eta\ll l'
\end{equation}
at every existing child scale.  The moment weight therefore changes
by at most a fixed factor; the same fixed inflation of
$m_1=C K^{n+1}m_0$ bounds both balls and moments.  Time shifts are
$o(l')$, so a fixed slightly enlarged native time range is preserved
through the bounded depth.

In particular, a depth-$j$ branch with $h=K^j$, $l=L/h$, has
\begin{equation}\label{s:node-parameter}
 m_j\le C^j h^{n+1}m,
\end{equation}
and the narrow prefactor, before the arbitrarily small decoupling
losses, is $C^jh^{-\beta}$.  The main factors cancel exactly:
\begin{equation}\label{s:cancellation}
 K^{2-(n+2)\beta}K^{(n+1)\beta}K^{-2n/(n+1)}=1,
 \qquad
 h^{-\beta}m_j^\beta l^{2n/(n+1)}
       \le C^j m^\beta L^{2n/(n+1)}.
\end{equation}

\subsection{Stopping and the terminal square-root-scale gain}

Stop just before a further division by $K$ would take the packet scale
below $D^{1/2}$.  Every existing scale has $l\ge D^{1/2}$, and terminal
scales satisfy
\begin{equation}\label{s:terminal-scale}
 D^{1/2}\le l<KD^{1/2}.
\end{equation}
The depth $J=O(v^{-1})$ is fixed once $v$ is chosen.

At any node work spatial box by spatial box at side $O(l^2)$.
Retain the exact packet labels whose original label trajectories meet a
fixed enlargement of the box.  Their energies have bounded overlap
over the boxes on the node time range, also for every child subarray.
The exact layers have arbitrarily prescribed decay about these
trajectories, uniformly after normalization.  Discarded fields are
therefore smaller than any chosen negative power of $L$, by
phase-space counting and Cauchy, since $l\ge L^{\tau/2}$.
Use the same retained labels for the total field and all its subfields.
Ignore polynomially tiny total-field levels before the activity tests,
and choose localization errors still smaller; this preserves the broad
lower bounds.  Equivalently apply the tests anew to the localized
exact array.  Its total and subfield input norms are bounded by retained
coefficient energy through Lemma~\ref{d:layers}.
These observations justify summing Lemma~\ref{s:broad} over spatial
boxes and summing energies over each depth of the tree.

At a terminal node, keep only $K^2$-cubes within transverse distance
$lL^\eta$ of at least one label trajectory; the others have negligible
contribution by the same rapid decay.  Restrict the node measure to
these cubes.  If a radius-$l$ ball meets this retained support, choose
one nearby trajectory from a cube it meets.  Since $K^2\ll l$ and
velocities are bounded, every point of the ball is at normalized
transverse distance $O(L^\eta)$ from that one trajectory.  Its moment
bound gives
\begin{equation}\label{s:terminal-mass}
 \nu(B(z,l))\le C L^{O(M\eta)}\frac{m_jl^{n+1}}D.
\end{equation}
The assertion also holds for fixed ball enlargements.  One trajectory
is used for each ball; there is no summation over nearby packets.

Normalize the node energy to one and sort retained cubes by
$\nu(Q)\asymp m_j\chi$, where $\chi\lesssim K^{2n}$.  Polynomially
tiny levels are negligible by the polynomial cell count and exact-input
supremum bound, leaving only $O(\log L)$ levels.  In each $Q$ choose a
unit witness $q_Q$ of largest field supremum.  Since
$|Q|^{2/P}=K^{2(n-1)}$,
\begin{equation}\label{s:terminal-witness}
 \sum_{Q\text{ in the class}}\nu(Q)^\beta\norm V_{L^P(Q)}^2
 \lesssim K^{2(n-1)}(m_j\chi)^\beta
                                  \sum_Q\sup_{q_Q}|V|^2.
\end{equation}
These witnesses need not carry measure.  Their counts are controlled
by their disjoint full cubes, each of mass comparable to $m_j\chi$.
For $r\ge1$ the count is at most
$C\chi^{-1}(r+CK^2)^n\lesssim K^{2n}\chi^{-1}r^n$.
At radius $l$, \eqref{s:terminal-mass} gives the improved count
$C L^{O(M\eta)}\chi^{-1}l^{n+1}/D$.

Apply Theorem~\ref{d:dz-input} at $S\asymp l^2$ with parameters
\[
 \gamma_0\lesssim K^{2n}\chi^{-1},\qquad
 \lambda_0\lesssim L^{O(M\eta)}\chi^{-1}l^{n+1}/D.
\]
Writing
\[
 a_n=\frac4{(n+1)(n+2)},\qquad
 \theta_n=\frac{2n}{(n+1)(n+2)},
\]
its bound is $O_\eps(S^\eps)\gamma_0^{a_n}\lambda_0^{\theta_n}
S^{\theta_n}$.  The normalized weight cancels exactly because
$\beta-a_n-\theta_n=0$.  Combining with \eqref{s:terminal-witness}
and restoring energy proves
\begin{equation}\label{s:terminal-bound}
 \mathcal T_\beta(V,\nu)
 \lesssim_\eps K^{C_n}L^{O(M\eta)+\eps}
       m_j^\beta l^{2n/(n+1)}(l/D)^{\theta_n}\mathcal E_{\rm node},
\end{equation}
where one may take
$C_n=2(n-1)+8n/((n+1)(n+2))$, increasing it for fixed grid
enlargements.  Indeed the unsimplified $l$ power is
$(n+3)\theta_n=2n/(n+1)+\theta_n$.
The logarithmic level count has been absorbed in the freely chosen
small power loss.  The input localization above handles spatial
supports larger than one $l^2$-box.

\subsection{Order of parameters and summation of the iteration}

\begin{proof}[Proof of Proposition~\ref{s:sparse}]
It remains to make the positive gains dominate every loss uniformly
in $D$.  Bounded $L$ is absorbed in the constant, and $\tau>1$ has no
unbounded admissible scales, so assume $0<\tau\le1$.
Fix first
\begin{equation}\label{s:saving-choice}
 a_* =\min\{c_b/2,\theta_n/2\},\qquad c_s=a_*/8.
\end{equation}
A broad leaf gains $l^{-c_b}\le D^{-c_b/2}$, and a terminal leaf
gains, by \eqref{s:terminal-scale},
\[
 (l/D)^{\theta_n}\le K^{\theta_n}D^{-\theta_n/2}.
\]
Thus every leaf gains $D^{-a_*}$ before its small losses and the
displayed terminal $K$ factors.  In particular this preliminary gain
is independent of $\tau$.

Choose $v>0$ sufficiently small, with fixed strict margins, that
\begin{equation}\label{s:parameter-inequalities}
 v<\frac{\tau c_b}{4},\qquad
 2v<\frac\tau4,\qquad
 (C_n+\theta_n)v<\frac{\tau a_*}{16}.
\end{equation}
These imply $K\le l^{c_b}$ at all broad nodes, and permit the terminal
$K$ factors to consume at most $D^{a_*/8}$ after fixed rounding
constants are absorbed.  The depth $J=O(v^{-1})$ is now fixed.
Choose the decoupling and refined-fractal losses so that their total
product along a path is at most $D^{a_*/8}$.  Choose $\eta>0$ next,
depending on $J,M$ as well, so that all padding losses, including
\eqref{s:terminal-mass}, cost at most $D^{a_*/8}$, and
\[
 2v+\eta<\tau/2.
\]
This proves \eqref{s:shift-smallness} at every existing child scale,
because $l'\ge D^{1/2}\ge L^{\tau/2}$.  Choose
$\gamma<\tau/8$, so that the retained layer shift satisfies
$|k|/l=o(1)$ at every node.  Finally choose the frequency-tail powers,
spatial decay orders, and profile differentiability large enough to
make all discarded errors negligible.  Those last choices do not
change $c_b$ or $c_s$.

On a depth-$j$ branch, \eqref{s:cancellation} makes the transported
quantity $m_j^\beta l^{2n/(n+1)}$ equal, up to $C^j$, to the root
quantity after multiplication by the narrow prefactor.  At a fixed
depth, angular bins partition energies, separate shifts have a
uniformly summable total weight, finite neighbor choices cost a fixed
constant, and spatial localizations have bounded energy overlap.
Thus the total weighted leaf energy at that depth is at most $C^j$
times the original energy.  Summing over the bounded number of depths
costs only a constant depending on $\tau$.
The allocations above leave more than the saving $D^{-c_s}$.

For clarity, ignored errors are also uniform after the whole tree is
summed.  On each original spatial box normalize mass parameter and
energy; there are polynomially many labels after localization, cells,
shift choices after truncation, and branches at bounded depth.  Their
parameters are polynomially bounded.  The tail orders were chosen last
and therefore beat this entire fixed polynomial, relative to
$m^\beta L^{2n/(n+1)}D^{-c_s}\mathcal E$.
Minkowski over the integrable exact-layer weights, the Sobolev reduction
\eqref{s:sobolev}, and the initial spatial energy sum now give
\eqref{s:estimate}, also for the prescribed derivatives.
\end{proof}
